\section{Automatic extraction of the characteristic function from a simulation model}\label{sec:extraction}

In engineering practice, a system studied for stability almost
always exists as a \emph{simulation model} first: a right-hand side
function written for a general-purpose integrator such as
\pkg{DifferentialEquations.jl}~\cite{rackauckas2017differentialequations}.
Deriving the characteristic function~\eqref{eq:charfun} from it by hand is
tedious and error-prone, particularly for many degrees of freedom, several
delays, or algebraic constraints. The package therefore extracts
$\Dfun(\lam)$ \emph{algorithmically} from the very same right-hand side
$\vect f(\vect u,\vect h,\vect p,t)$ used for time-domain simulation, where
$\vect h(\vect p, t-\tau)$ is the standard history-function interface for
the delayed states.

The mechanism requires a single evaluation of $\vect f$ per requested
$\lam$. The state is seeded with dual-number perturbations,
$\vect u = \vect u_{\mathrm{eq}} + \vect a$, where the components of
$\vect a$ carry independent infinitesimal directions, and the history
argument is replaced by the \emph{mock exponential history}
\begin{equation}\label{eq:mockhistory}
  \vect h(\vect p, t) = \vect u_{\mathrm{eq}} + \vect a\,\ee^{\lam t}.
\end{equation}
Whatever delayed value $\vect h(\vect p, t-\tau_j)$ the model requests
contributes the factor $\ee^{-\lam\tau_j}$ automatically, so discrete,
multiple, and (approximated) distributed delays are all captured without
declaring them. Evaluating $\vect f$ at $t=0$ and collecting the dual
partial derivatives of the residual
\begin{equation}\label{eq:extraction}
  \vect r = \lam\,\mat E\,\vect a - \vect f\bigl(\vect u_{\mathrm{eq}}
      + \vect a,\ \vect h,\ \vect p,\ 0\bigr)
  \qquad\Longrightarrow\qquad
  \mat M(\lam) = \frac{\partial\vect r}{\partial\vect a},\quad
  \Dfun(\lam)=\det\mat M(\lam),
\end{equation}
yields the exact characteristic matrix of the \emph{linearized} dynamics
around $\vect u_{\mathrm{eq}}$: the dual seeding performs the linearization,
so the right-hand side may be fully \emph{nonlinear} -- no separate
linearization step is required (the equilibrium itself could also be located
from the same right-hand side, which is outside the scope of this paper; a
residual check warns if the supplied $\vect u_{\mathrm{eq}}$ is not an
equilibrium). The mass matrix $\mat E$ is supplied separately -- it is not
part of the right-hand side, exactly as in the
\code{ODEFunction}/\code{DDEFunction} convention of
\pkg{DifferentialEquations.jl} -- and it may be \emph{singular}: rows of
zeros in $\mat E$ turn the corresponding rows of $\mat M(\lam)$ into the
algebraic constraint equations, so delay differential--algebraic systems are
handled by the identical code path. The determinant is evaluated with
statically sized arrays, so a full extraction costs one right-hand-side call
plus an $N\times N$ determinant, and the dual seeding composes transparently
with the outer dual numbers of the phase ODE~\eqref{eq:phaseode}.

The complete user workflow is thus reduced to providing the simulation
right-hand side, the mass matrix if it differs from the identity, and the
parameter ranges of the chart:
\begin{verbatim}
function rhs(u, h, p, t)              # simulation model (DDE interface)
    P, D = p
    x1_del = h(p, t - 0.5; idxs=1)    # delayed position
    v1_del = h(p, t - 0.5; idxs=4)    # delayed velocity
    ...                               # equations of motion + constraints
end

D(lam, p)  = get_D_from_model(rhs, lam, p, Val(7); mass_matrix = E)
Z, Zraw, minD, sigma_est, omega_est =
    calculate_unstable_roots_direct(D, (2.0, 1.0))
\end{verbatim}
The leading order $n$ required by the counting
formula~\eqref{eq:count} is estimated automatically
by~\eqref{eq:npow} -- including the effective orders of neutral
characteristic functions and the non-integer orders of fractional ones --
so no additional structural input is needed. In
Section~\ref{sec:showcase} the extraction is verified against a hand-derived
characteristic function: the two agree to machine precision.